\subsection{中心单代数}

定义 1. —— 若映射 \(\lambda\mapsto\lambda1\) 是从 K 到 A 的中心的一个双射，则称域 K 上的代数 A 是中心的。

中心代数不化为 0。对每个整数 \(n \geqslant 1\)，矩阵 K-代数 \(\mathbf{M}_{n}(\mathrm{K})\) 是中心的（VIII, 第 83 页, 推论 2）且是单的（VIII, 第 120 页, 定理 1）。更一般地，设 D 是一个有限次数的中心 K-代数；则 \(\mathbf{M}_{n}(\mathrm{D})\) 也是中心的。设 A 是一个单环；它的中心 Z 是一个域（VIII, 第 121 页, 推论 1），因此 A 是 Z 上的一个中心单代数。若域 K 是代数闭的，则 K 上的一个有限次数的单代数是中心的（VIII, 第 122 页, 推论 3）。中心单代数的反代数是中心单的。

注. —— 1) 设 A 与 B 是 K-代数。若 \(A \otimes_{K} B\) 是一个中心单代数，则 A 与 B 亦然（III, §4, No.~4, 第 468 页, 命题 6 的推论与 VIII, 第 221 页, 命题 6）。反之，若 A 与 B 都是中心单代数且其中之一在 K 上次数有限，则 \(A \otimes_{K} B\) 是一个中心单代数（III, §4, No.~4, 第 468 页, 命题 6 的推论与 VIII, 第 222 页, 推论 2）。

\begin{enumerate}
\def\labelenumi{\arabic{enumi})}
\setcounter{enumi}{1}
\item
  设 A 是一个 K-代数，L 是域 K 的一个扩张。若 L-代数 \(A_{(L)}\) 是中心单的，则 K-代数 A 是中心单的。反之，若次数 {[}A : K{]} 与 {[}L : K{]} 之一有限且 A 是一个中心单 K-代数，则 L-代数 \(A_{(L)}\) 是中心单的。这由 VIII, 第 222 页的推论 2 得出。
\item
  设 A 与 B 是森田等价的 K-代数。代数 A 是中心单的当且仅当 B 是中心单的（VIII, 第 105 页, 命题 6；第 111 页, 推论；以及第 102 页, 推论 1）。
\item
  特别地，若 A 是一个中心单 K-代数且 \(n \geqslant 1\)，则 \(\mathbf{M}_{n}(\mathrm{A})\) 是一个中心单 K-代数（VIII, 第 102 页, 例 1）。
\end{enumerate}

定理 1. —— 设 A 是一个有限次数的 K-代数。下述性质等价：

\begin{enumerate}
\def\labelenumi{(\roman{enumi})}
\item
  代数 A 是中心单的。
\item
  代数 \(\mathbf{A}\) 是中心的且无根。
\item
  从 K-代数 \(A \otimes_{K} A^{\circ}\) 到 K-代数 \(\operatorname{End}_{K}(A)\) 的把 \(a \otimes a'\) 映到从 A 到 A 的 K-线性映射 \(x \mapsto axa'\) 的典范同态是双射的。
\item
  存在域 K 的一个扩张 L 与一个整数 \(n \geqslant 1\)，使得 L-代数 \(\mathrm{A}_{(\mathrm{L})}\) 与 \(\mathbf{M}_{n}(\mathrm{L})\) 同构。
\item
  对每个可分闭包 \(\mathrm{K}'\)（\(\mathbf{K}\) 的），存在一个整数 \(n \geqslant 1\)，使得 \(\mathrm{K}'\) -代数 \(\mathrm{A}_{(\mathrm{K}')}\) 与 \(\mathbf{M}_n(\mathrm{K}')\) 同构。
\item
  存在域 K 的一个有限次数的伽罗瓦扩张 L 与一个整数 \(n \geqslant 1\)，使得 L-代数 \(\mathrm{A}_{(\mathrm{L})}\) 与 \(\mathbf{M}_n(\mathrm{L})\) 同构。
\item
  存在一个有限次数的 K-代数 D（它是一个以 K 为中心的域）与一个整数 \(n \geqslant 1\)，使得代数 A 同构于代数 \(\mathbf{M}_{n}(\mathrm{D})\)。
\end{enumerate}

一个环是单的当且仅当它是半单的且它的中心是一个域（VIII, 第 143 页, 命题 10 的推论）。由于 A 是域 K 上的一个有限次数代数，它是一个左阿廷环；因此它半单当且仅当它无根（VIII, 第 154 页, 命题 4）。由此得出 (i) 与 (ii) 的等价性。

令 \(E = A \otimes_{K} A^{\circ}\) 与 \(F = \text{End}_{K}(A)\)；用 \(\varphi\) 表示由关系式 \(\varphi(a \otimes a')(x) = axa'\)（对 A 中的 \(x, a, a'\)）定义的从 E 到 F 的典范同态。若代数 A 是中心单的，则 \(A^{\circ}\) 亦然，从而 E 亦然（注 1），于是 \(\varphi\) 是单射的。现在我们有 \([E : K] = [A : K]^{2} = [F : K]\)，故 \(\varphi\) 是双射的。反之，假设 \(\varphi\) 是双射的；由于代数 F 是中心单的（因为它同构于一个矩阵代数 \(\mathbf{M}_{m}(K)\)），E 亦然，从而 A 亦然（注 1）。我们这样就证明了 (i) 与 (iii) 的等价性。

由注 4，断言 (vii) 蕴涵断言 (i)。其逆蕴涵由 VIII, 第 122 页的推论 3 与 VIII, 第 83 页的推论 2 得出。

显然 (vi) 蕴涵 (iv)，而由注 2，(iv) 蕴涵 (i)。

剩下要证明蕴涵关系 (i)⇒(v)⇒(vi)。假设 A 是中心单的，并设 \(K'\) 是 K 的一个可分闭包（V, §7, No.~8, 第 45 页）。则 \(\mathrm{A}_{(\mathrm{K}')}\) 是 \(K'\) 上的一个有限次数的中心单代数（VIII, 第 251 页, 注 2）。于是由 VIII, 第 231 页的推论，存在一个整数 \(n \geqslant 1\) 与一个 \(K'\) -代数同构（从 \(\mathrm{A}_{(\mathrm{K}')}\) 到 \(\mathbf{M}_{n}(\mathrm{K}')\)）；注意 \(K'\) -代数 \(\mathbf{M}_{n}(\mathrm{K}')\) 与 \(\mathbf{M}_{n}(\mathrm{K})_{(\mathrm{K}')}\) 同构。由 VIII, 第 234 页的引理 4，存在 \(K'\) 的一个在 K 上有限生成的子扩张 L，使得 L-代数 \(\mathrm{A}_{(\mathrm{L})}\) 与 \(\mathbf{M}_{n}(\mathrm{K})_{(\mathrm{L})}\) 同构。于是 L 在 K 上是可分的且次数有限，故包含在一个在 K 上伽罗瓦且次数有限的子扩张 \(L'\)（\(K'\) 的）中（V, §10, No.~1, 第 57 页, 命题 2）。于是 \(L'\) -代数 \(\mathrm{A}_{(\mathrm{L}')}\) 与 \(\mathbf{M}_{n}(\mathrm{L}')\) 同构。

推论 1. —— 设 A 是可分闭域 K 上的一个有限次数的中心单代数。存在一个整数 \(n \geqslant 1\)，使得 A 同构于矩阵代数 \(\mathbf{M}_{n}(\mathrm{K})\)。

实际上，\(\mathrm{K}\) 的每个伽罗瓦扩张都等于 \(\mathrm{K}\)；只需应用定理 1 的性质 (i) 与 (v) 的等价性。

推论 2. —— 设 A 是 K 上的一个有限次数的中心单代数（例如，K 上的一个以 K 为中心的有限次数的域）。存在一个整数 \(n \geqslant 1\)，使得 \([A : K] = n^{2}\)。

设 \(\mathrm{L}\) 是 \(\mathrm{K}\) 的一个扩张，\(n\) 是一个严格正的整数，使得 L-代数 \(\mathrm{A}_{(\mathrm{L})}\) 与 \(\mathbf{M}_n(\mathrm{L})\) 同构。我们有

\[
[ \mathrm{A}: \mathrm{K} ] = [ \mathrm{A} _ {(\mathrm{L})}: \mathrm{L} ] = [ \mathbf {M} _ {n} (\mathrm{L}): \mathrm{L} ] = n ^ {2}.
\]

在推论 2 的记号中，整数 n 称为 A 的约化次数。

注 5. —— 设 A 是 K 上的一个约化次数为素数 \(\ell\) 的有限次数中心单代数。则 A 或者是一个域，或者同构于 \(\mathbf{M}_{\ell}(\mathrm{K})\)。实际上，A 同构于形如 \(\mathbf{M}_{n}(\mathrm{D})\) 的一个代数，其中 D 是一个以 K 为中心的域，且我们有

\[
\ell^ {2} = [ \mathrm{A}: \mathrm{K} ] = n ^ {2} [ \mathrm{D}: \mathrm{K} ];
\]

若 A 不是一个域，则 \(n \neq 1\)，故 \(n = \ell\) 且 D = K。

